\section{Foundation providers and parameter choices}\label{sec:interface}

Section~\ref{sec:inputs} states the complete hypotheses used in this
article.  Their proofs are in the separate foundation companion
\emph{Uniform excess charge for Coulomb molecules and the outer radius
of neutral atoms} \cite{OuterRadius}.  This appendix identifies the
provider results and explains the compatibility of their parameter
choices.  All inputs use two spin states, kinetic operator $-\Delta/2$,
and the constants in \eqref{intro:constants}.

\paragraph{States and raw counts.}
Ground-state existence comes from \FoundationBinding\ of
\cite{OuterRadius}, and sector monotonicity from \FoundationMonotonicity;
the bounded neutral offset uses \FoundationOffset.  The annular $L^2$ bound in
\cref{inp:energy-counts} is \FoundationCounts, and the unit-radius
exterior bound follows from \FoundationDeficit\ with $t=4$ and exactly
$Z$ electrons.  These inputs apply uniformly to every normalized neutral
ground state, not only to a selected family of states.

\paragraph{Common observations and conditional versions.}
\FoundationObservations\ of \cite{OuterRadius} supply the observation
law and likelihood-integral versions used in
\eqref{inp:observations}--\eqref{inp:mass-poisson}.  In particular, the
same master packets define all the conditional densities.  Their exact
mass, Poisson and tower identities and finite-system regularity hold
on fixed full-measure sets.  The derivative and local Lipschitz bounds
need not be uniform in $Z$.

\paragraph{Comparison and intermediate subsolutions.}
\FoundationComparison\ of \cite{OuterRadius} supplies
\cref{inp:comparison-prop}: arbitrary fixed annular width and accuracy,
a universal upper cap, and an exceptional probability of order
$r_j^{25}$.  \FoundationIntermediate\ supplies
\cref{inp:barrier-prop} at every intermediate index, with constants
independent of the number of scales.  Its expected error bound controls
every prefix, as explained after \eqref{inp:barriers}.  At the terminal
index the sigma-field is trivial, so the same bound is deterministic.
The almost-sure subsolution assertions do not make the comparison event
a full-measure event.

\paragraph{Order of choices.}
First fix the kernel and propagation constants, including $\eps$ and
the positive barrier coefficient $B$.  Then fix a requested comparison
band and accuracy, make $s$ sufficiently small for that comparison and
the fixed barrier construction, and take $Z$ above the finite threshold
$Z_{\mathrm{bar}}(s)$.  This threshold includes $Z\ge Z_0$ and
$2\eps Z^{-1/3}\le s$.  Changing $L_1,h_l,h_h,\xi$ changes the comparison
event and its smallness requirement, not the observations, packets,
conditional fields, or fixed barrier construction.

These two uniformities serve different purposes.  The common fields
allow Section~\ref{sec:scales} to improve the comparison accuracy along
a slow diagonal while retaining the same subsolutions.  Uniformity in
the neutral ground state allows Section~\ref{sec:tails} to turn the
result along every admissible sequence into a finite-charge threshold
valid for all such states.  No rate of growth for
$Z_{\mathrm{bar}}(s)$ is asserted or used.
